% Authored: observability.

\chapter{Observability}
\label{ch:observability}

\section{Two layers}

\cref{module:feedback_intelligence_agent.telemetry} provides two independent
mechanisms, and the distinction matters operationally.

\begin{description}
  \item[Structured logging] \texttt{configure\_logging} sets the root logger to
    emit bare messages, and \texttt{log\_event(name, payload)} writes a single
    JSON object per event with sorted keys. Log events are always emitted; they
    are not gated by the telemetry switch. They are what a container's stdout
    collector sees.
  \item[Span telemetry] \texttt{Telemetry} records events and spans with
    timestamps, durations, and a correlation identifier, through a pluggable
    sink. A \texttt{Telemetry} constructed without a sink is a no-op: the
    default code path emits nothing and pays nothing.
\end{description}

\section{Sinks}

Three sinks implement the \texttt{TelemetrySink} protocol:

\begin{description}[nosep]
  \item[\texttt{InMemoryTelemetrySink}] collects events for assertions; used by
    the test suite.
  \item[\texttt{JsonlTelemetrySink}] appends one JSON object per line to a
    configured file, guarded by a lock so concurrent emissions do not
    interleave. This is the default when telemetry is enabled.
  \item[\texttt{OpenTelemetryTraceSink}] converts events into spans on the
    process-global OpenTelemetry tracer provider. It requires the \texttt{otel}
    extra and raises \texttt{TelemetryConfigurationError} when the SDK is
    unavailable, rather than silently degrading.
\end{description}

\section{Instrumented stages}

Spans are emitted with a shared correlation identifier per request, so a single
answer can be reconstructed end to end:

\begin{longtable}{p{0.32\textwidth} p{0.58\textwidth}}
\toprule
Span & Recorded attributes \\
\midrule
\endfirsthead
\toprule
Span & Recorded attributes \\
\midrule
\endhead
\bottomrule
\endlastfoot
\texttt{ingestion\_*} & path, strict mode, record count \\
\texttt{embedding\_*} & model, dimension, chunk count, vector count \\
\texttt{agent\_run\_*} & route, top-$k$, question length, guardrail decision,
  dropped chunks, tool name and status, citations, confidence, retrieved
  chunks, hallucination risk and score \\
\texttt{retrieval\_*} & retriever class, route, top-$k$, candidate-$k$,
  metadata filters, filtered-out count, reranker class, result count,
  min and max score \\
\texttt{tool\_run\_*} & tool name, selection reason, output field count \\
\texttt{llm\_call\_*} & provider, prompt characters, response characters \\
\texttt{response\_parsing\_*} & output format, repair applied, validation error \\
\texttt{hallucination\_check\_*} & risk, detection flag, overlap score,
  unsupported sentence count, judge label and confidence \\
\texttt{evaluation\_finished} & every aggregate metric of the report \\
\end{longtable}

Notice what these attributes are: they are the quality signals of the system,
not just timings. Retrieval score spread, citation count, hallucination risk,
repair rate, and guardrail drops are all observable per request without
re-running anything.

\section{Health and readiness}

\texttt{GET /health} is a liveness probe: it performs no dependency check, so
an orchestrator can use it to decide whether to restart the container.
\texttt{GET /ready} is a readiness probe whose semantics follow from the
startup design --- because the agent and all stores are constructed during
application construction, reaching the handler proves the application is fully
built. Neither endpoint is rate limited and neither requires a key.

The production Compose file wires a container healthcheck; the ECS task
definition and the Fly.io configuration are the other deployment descriptors
present.

\section{What is not implemented}

There are no metrics counters, no histogram export, no log aggregation
configuration, no dashboards, and no alerting rules in the repository. The
telemetry attributes above are the raw material for such things, not the things
themselves.
